\documentclass[a4paper,12pt]{report}

\usepackage{geometry}
\geometry{margin=1in}
\usepackage{graphicx}
\usepackage{amsmath}
\usepackage{hyperref}
\usepackage{listings}
\usepackage{xcolor}
\usepackage{tabularx}
\usepackage{enumitem}
\setlist[itemize]{itemsep=0.5em,leftmargin=2em}




\lstset{%
    basicstyle=\ttfamily\small,
    frame=single,
    backgroundcolor=\color{gray!10},
    keywordstyle=\color{blue},
    commentstyle=\color{green!50!black},
    stringstyle=\color{orange},
    numbers=left,
    numberstyle=\tiny,
    numbersep=5pt,
    showstringspaces=false,
    breaklines=true
}

\begin{document}

\begin{titlepage}
    \begin{minipage}[t]{0.5\textwidth}
        \includegraphics[width=1\textwidth]{ASP.jpeg}
    \end{minipage}
    \hfill
    \begin{minipage}[t]{0.5\textwidth}
        \raggedleft
        \includegraphics[width=1\textwidth]{IAP.jpg}
    \end{minipage}


    \vspace{1cm}

    {\Huge\centering\bfseries Internship Report\\[0.5cm]}
    \vspace{1cm}

    \raggedright
    \large

    {\centering\bfseries Topic: Report on Chromatix\\}

    \vspace{2cm}

    \centering\textbf{Author:}\\[0.5cm]
    \centering Ziyi Xiong\\[0.5cm]
    \centering Master Student in Photonics at ASP\\[0.5cm]
    
    \centering\textbf{Internship Period:}\\[0.5cm] 
    \centering  01-09-2024 to 28-02-2025\\[0.5cm]
    
    \centering \textbf{Supervisor:} \\[0.5cm]
    \centering Prof. Dr. Vladan Blahnik, Institute of Applied Physics\\[0.5cm]
    
    \centering\textbf{Industry Advisor:} \\[0.5cm]
    \centering Prof. Dr. Frank Wyrowski, President of LightTrans GmbH\\

    \vfill
\end{titlepage}


\tableofcontents
\newpage

\chapter{Introduction}
\section{Context of the Internship}
The internship was conducted under the joint guidance of my university supervisor, Prof. Dr. Vladan Blahnik, and my industry advisor, Prof. Dr. Frank Wyrowski. The project entailed a deep-dive analysis into the open-source Python library known as \textbf{Chromatix}, which was originally developed at HHMI Janelia Research Campus. This library leverages wave-optics and advanced computational frameworks (notably JAX and FLAX) to enable \textit{differentiable} simulations of optical systems.\newline
\newline
My role was to thoroughly understand Chromatix, from its underlying wave-optical theory to the structure of its Python source code, and to explore real use cases. Specifically, I examined how classes, functions, and modules contribute to building optical elements, performing wave propagation, and enabling inverse-design tasks such as optimizing lens parameters and masks.

\section{Goals and Objectives}
The primary goals of this internship were:
\begin{enumerate}
    \item \textbf{Comprehensive Code Analysis:} Understand and document the organization and functionality of the Chromatix library (classes, functions, modules).
    \item \textbf{Theoretical Foundations:} Relate each function's or class's implementation to the fundamental physics of wave optics (e.g., Fresnel diffraction, Fourier transforms, Jones calculus, etc.).
    \item \textbf{Practical Demonstrations:} Demonstrate usage through multiple example systems (simple widefield PSF simulation, a 4f system, and a CGH setup using a Digital Micromirror Device (DMD)).
    \item \textbf{Evaluation \& Future Prospects:} Critically assess strengths, limitations, and propose future directions for library development.
\end{enumerate}

\section{Structure of This Report}
This report is divided into the following chapters:
\begin{itemize}
    \item \textbf{Chapter 1: Introduction} - Provides an overview of the internship objectives and a brief introduction to Chromatix. Delves into the design principles, module structure, and theoretical basis of the library.
    \item \textbf{Chapter 2: Core Source Code Analysis} - Explores the main classes and functions in Chromatix, mapping them to wave-optical principles.
    \item \textbf{Chapter 3: Practical Usage Examples} - Demonstrates real-world simulations using Chromatix with extended code snippets.
    \item \textbf{Chapter 4: Conclusions \& Future Work} - Summarizes key findings and offers recommendations for further improvements.
\end{itemize}

\chapter{Chromatix in Detail}
\section{Overview}
Chromatix is a differentiable wave-optics library designed to model complex optical systems while allowing gradient-based optimization, thanks to the JAX ecosystem. Key features include:
\begin{itemize}
    \item \textbf{GPU Acceleration}: Offloading large-scale computations to GPUs or TPUs.
    \item \textbf{Inverse Design}: Using auto-differentiation to iteratively refine system parameters (e.g., lens curvature, phase masks).
    \item \textbf{Modular Architecture}: Dividing code into multiple packages (\texttt{data}, \texttt{functional}, \texttt{elements}, etc.) for clarity and maintainability.
    \item \textbf{Multi-Wavelength \& Polarization Handling}: Scalar and vector fields, and spectral weighting.
\end{itemize}

\section{Package Breakdown}
\textbf{Chromatix} is organized into six major packages plus two additional modules:
\begin{itemize}
    \item \textbf{data}: Generates fundamental optical patterns (radial shapes, grayscale images) and can create permittivity tensors for materials. While not central to wave-propagation, it provides ready-made data for simulation.
    \item \textbf{functional}: Implements physical transformations like Fresnel or Angular Spectrum propagation, lens-induced phase changes, amplitude/phase masks, polarizers, and sensor operations. This package is the mathematical backbone, reliant heavily on JAX (\texttt{jax.numpy}).
    \item \textbf{elements}: Encapsulates optical components as trainable modules (subclasses of \texttt{flax.linen.Module}). For example, \textit{lens elements}, \textit{phase mask elements}, etc. This design allows each element to be optimized via gradient-based methods.
    \item \textbf{systems}: Provides the \texttt{OpticalSystem} class to combine multiple \textit{elements} into a pipeline. Example systems include a standard microscope (\texttt{Microscopy}) and a 4f system (\texttt{Optical4FSystemPSF}).
    \item \textbf{ops}: Supplies typical image-processing utilities (filters, noise, quantization) that can be integrated into pipeline simulations.
    \item \textbf{utils}: Hosts a variety of helper functions (FFT wrappers, 2D/3D shape generation, Zernike polynomials, wavefront aberrations, etc.) to support the core simulation tasks.
    \item \textbf{field.py}: Defines the \texttt{Field} class (with \texttt{ScalarField} and \texttt{VectorField}) that serve as the data structure for complex wave fields. Includes sampling intervals, wavelength arrays, and advanced indexing.
    \item \textbf{\_\_init\_\_\.py}: Basic initializations, ensuring that the library is well-structured upon import.
\end{itemize}

\section{Core Theoretical Concepts}
Chromatix leverages fundamental wave optics:
\begin{enumerate}
    \item \textbf{Fourier Optics}: The library heavily uses Fourier transforms to handle propagations (Fresnel, Fraunhofer, Angular Spectrum).
    \item \textbf{Jones Calculus}: Particularly for vector fields, polarizers, and wave plates.
    \item \textbf{Sampling Theory}: Ensuring that the discretized fields respect sampling requirements, particularly for FFT-based methods.
    \item \textbf{Differentiable Programming}: JAX's automatic differentiation engine allows every wave transform to be differentiable w.r.t. design parameters.
\end{enumerate}

\chapter{Core Source Code Analysis}
\section{Field Class (\texttt{field.py})}
The \textbf{Field} class (and its scalar/vector variants) is the center piece:
\begin{itemize}
    \item \textbf{u}: A complex array representing the electromagnetic field samples. Typically shaped as $(B..., H, W, C, [1|3])$, where $H\times W$ is the spatial grid, $C$ is the number of spectral channels, and $[1|3]$ indicates scalar or vector fields.
    \item \textbf{\_dx}: Sampling intervals (y, x) specifying real-space pixel sizes.
    \item \textbf{\_spectrum}: One or more wavelengths used in the field (multi-wavelength simulations supported).
    \item \textbf{\_spectral\_density}: Respective weights of each wavelength.
\end{itemize}
\noindent \textbf{Key Methods} include:
\begin{itemize}
    \item \texttt{grid()}: Returns spatial coordinate arrays, centered in the middle of the array.
    \item \texttt{k\_grid()}: Returns frequency coordinates (Fourier domain) suitable for propagation computations.
    \item \texttt{phase()}: Extracts the phase of the complex field.
    \item \texttt{intensity()}: Computes the intensity $I=|u|^2$.
    \item \texttt{amplitude()}: Computes the amplitude $|u|$.
    \item Arithmetic operators (\texttt{\_\_add\_\_}, \texttt{\_\_mul\_\_}, etc.) that permit easy combinations of fields.
\end{itemize}

\noindent\textbf{Physical Rationale}: By bundling sampling intervals and spectral data within \texttt{Field}, Chromatix simplifies transform-based operations (Fourier transforms, lens phases, etc.). The methods ensure consistent coordinate handling.

\subsection{Why JAX?}
JAX transforms each function into a computation graph that can be differentiated. For wave optics, each transformation (e.g., lens phase or Fresnel propagation) is re-cast in a way that JAX can keep track of partial derivatives. This is extremely useful for tasks such as inverse design, where we want to adjust lens surfaces or mask patterns to achieve a target field.

\section{Functional - Module \& Function Overview}
The \textbf{functional} package provides static (pure) functions implementing standard optical transformations:

% 1. amplitude_masks.py
\paragraph*{1. \textbf{amplitude\_masks.py}}
\begin{itemize}
  \item \texttt{amplitude\_change}: This function perturbs a Field object by a given amplitude array.
\end{itemize}

% 2. convenience.py
\paragraph*{2. \textbf{convenience.py}}
\begin{itemize}
  \item \texttt{optical\_fft}: This \texttt{optical\_fft} function performs a Fresnel propagation on a Field object using a Fast Fourier Transform (FFT) or Inverse Fast Fourier Transform (IFFT) depending on the propagation direction.
\end{itemize}

% 3. lenses.py
\paragraph*{3. \textbf{lenses.py}}
\begin{itemize}
  \item \texttt{thin\_lens}: This function is to calculate the field after passing a thin lens.
  \item \texttt{ff\_lens}: This function simulates a situation where a field propagates a distance f to a thin lens, then it passes through the thin lens and then propagates another distance f to the focal plane of the lens.
  \item \texttt{df\_lens}: This function simulates a more general case of 2f setup, where the first f becomes d meaning that the initial field starts from any distance before the thin lens.
\end{itemize}

% 4. phase_masks.py
\paragraph*{4. \textbf{phase\_masks.py}}
\begin{itemize}
  \item \texttt{spectrally\_modulate\_phase}: This function ensures accurate modelling of optical systems involving multiple wavelengths.
  \item \texttt{phase\_change}: This function aims to perturb the input field by an input phase array.
  \item \texttt{wrap\_phase}: The function serves to adjust phase values so that they fall within a specified range.
\end{itemize}

% 5. polarizers.py
\paragraph*{5. \textbf{polarizers.py}}
\begin{itemize}
  \item \texttt{jones\_vector}: This function creates a Jones vector with the inputs theta and beta, where theta is the polarization angle and beta is the phase delay.
  \item \texttt{linear}: This function is for generating a Jones vector for linearly polarized light.
  \item \texttt{left\_circular}: This function generates a Jones vector for left-hand circularly polarized light.
  \item \texttt{right\_circular}: This function generates a Jones vector for right-hand circularly polarized light.
  \item \texttt{polarizer}: This function is deigned to simulate polarizers such as linear polarizer, left circular polarizer, and right circular polarizer.
  \item \texttt{linear\_polarizer}: This function applies a linear polarizer to the input field at a given angle with respect to the horizontal.
  \item \texttt{left\_circular\_polarizer}: This function applies a left circular polarizer to the incoming field.
  \item \texttt{right\_circular\_polarizer}: This function applies a right circular polarizer to the incoming field.
  \item \texttt{phase\_retarder}: This simulates a phase retarder by applying a general phase retarder Jones matrix to the input field to modify its polarization state.
  \item \texttt{wave\_plate}: This function simulates a general wave plate, also known as linear phase retarder by calling the phase\_retarder function.
  \item \texttt{halfwave\_plate}: Halfwave plate makes the light's orthogonal components have a phase difference of $\pi$ (half a wave) after the light passer through the phase retarder.
  \item \texttt{quarterwave\_plate}: Quarter wave plate makes the light's orthogonal components have a phase difference of $\pi/2$.
  \item \texttt{universal\_compensator}: A universal compensator, also known as universal polarizer, can be used to generate any polarization.
\end{itemize}

% 6. propagation.py
\paragraph*{6. \textbf{propagation.py}}
\begin{itemize}
  \item \texttt{transform\_propagate}: This function performs Fresnel propagation of an optical field over a specific distance using the Single-FFT Fresnel Propagation Method (SFT-FR).
  \item \texttt{compute\_sas\_precompensation}: This is to compute the precompensation factor used in the Scaled Angular Spectrum (SAS) method.
  \item \texttt{kernel\_propagate}: This function simulates to propagate a field using a given propagator with Fourier convolution.
  \item \texttt{transform\_propagate\_sas}: This function offers a method to do the single-FFT Fresnel propagation method with the precompensation factor.
  \item \texttt{compute\_transfer\_propagator}: This function computes the propagation kernel (transfer function) for the Fresnel propagation.
  \item \texttt{transfer\_propagate}: This function simply calls the function compute\_transfer\_porpgator to compute the propagator and the function kernel\_propgate to compute the propagated field.
  \item \texttt{compute\_exact\_propagator}: Compute propagation kernel for propagation with no Fresnel approximation.
  \item \texttt{exact\_propagate}: This function simply calls the function compute\_exact\_porpgator to compute the propagator and the function kernel\_propgate to compute the propagated field.
  \item \texttt{compute\_asm\_propagator}: This propagation kernel is also without Fresnel approximation, but the phase of the kernel could contain a shift in output plane.
  \item \texttt{asm\_propagate}: Propagate field for a distance z using angular spectrum method.
  \item \texttt{compute\_padding\_transform}: This is to compute the padding for transform propagation.
  \item \texttt{compute\_padding\_transfer}: Automatically compute the padding required for transfer propagation.
  \item \texttt{compute\_padding\_exact}: Automatically compute the padding required for exact propagation.
\end{itemize}

% 7. pupils.py
\paragraph*{7. \textbf{pupils.py}}
\begin{itemize}
  \item \texttt{circular\_pupil}: This function creates a circular mask and applies it to the incoming field.
  \item \texttt{square\_pupil}: This function creates a squared mask similarly to the circular\_pupil function.
\end{itemize}

% 8. samples.py
\paragraph*{8. \textbf{samples.py}}
\begin{itemize}
  \item \texttt{jones\_sample}: The jones\_sample function simulates how an incoming vector field is perturbed when it passes through a thin sample.
  \item \texttt{thin\_sample}: This is to simulate the perturbation to an incoming scalar field based on the thin sample approximation.
  \item \texttt{multislice\_thick\_sample}: The purpose of this function is to simulate a scalar field propagating through a thick sample by modelling the thick sample as multiple thin samples.
  \item \texttt{fluorescent\_multislice\_thick\_sample}: This function simulates the propagation of a scalar field through a thick, transparent, and fluorescent sample.
  \item \texttt{PTFT}: This function is to create a projection tensor for transverse wave (PTFT). This projection tensor projects vectors onto the plane perpendicular to the wave vector $\mathbf{k}$.
  \item \texttt{thick\_sample\_vector}: This function simulates a vectorial wave field with scattering potential propagating through a thick sample.
\end{itemize}

% 9. sensors.py
\paragraph*{9. \textbf{sensors.py}}
\begin{itemize}
  \item \texttt{basic\_sensor}: Produces an intensity image from an incoming field or intensity array and simulates shot noise.
\end{itemize}

% 10. sources.py
\paragraph*{10. \textbf{sources.py}}
\begin{itemize}
  \item \texttt{point\_source}: This function can generate a scalar field or vectorial field located at a distance z from the point source.
  \item \texttt{objective\_point\_source}: This function simulates the field that just after passing through a lens with focal length f and numerical aperture NA. The field is originated from the position which is deviated a distance z from the object space's focal plane.
  \item \texttt{plane\_wave}: This is to create field simulating a plane wave.
  \item \texttt{generic\_field}: The function creates either a scalar field or vectorial field with a self-defined phase which determines the shape of the wavefront.
\end{itemize}

\section{Elements Package}
The \textbf{elements} package wraps functions into \texttt{flax.linen.Module} classes, making them trainable modules.
% 1. amplitude_masks.py
\paragraph*{1. \textbf{amplitude\_masks.py}}
\begin{itemize}
  \item \texttt{AmplitudeMask}: Applies an amplitude mask to an incoming Field by calling the function amplitude\_change from the functional package. The amplitude can be learned pixel by pixel.
\end{itemize}

% 2. convenience.py
\paragraph*{2. \textbf{convenience.py}}
\begin{itemize}
  \item \texttt{Flip}: This element flips the incoming field upside down using jnp.flip().
  \item \texttt{ScaleAndBias}: This class simply adds a bias to the incoming field and multiplies the sum by a scale.
  \item \texttt{Binarize}: This class aims to binarize the incoming `Field` by calling the function binarize from chromatix.ops.quantization.
  \item \texttt{Quantize}: This class applies the quantize function from chromatix.ops.quantization to the incoming Field to quantize the data to a specific bit depth. 
\end{itemize}

% 3. lenses.py
\paragraph*{3. \textbf{lenses.py}}
\begin{itemize}
  \item \texttt{ThinLens}: It models the behavior of a thin lens in an optical system by taking three attributes (can be learned) to apply the effect of a thin lens to an incoming field and returns the modified field.
  \item \texttt{FFLens}:  Applies a thin lens placed a distance f after the incoming field. This element returns the field a distance f after the lens. The attributes f, n, and NA can be learned.
  \item \texttt{DFLens}: Applies a thin lens placed a distance d after the incoming field. This element returns the field a distance f after the lens. The attributes d, f, n, and NA can be learned.
\end{itemize}

% 4. phase_masks.py
\paragraph*{4. \textbf{phase\_masks.py}}
\begin{itemize}
  \item \texttt{PhaseMask}: Applies a phase mask to an incoming field by calling the phase\_change function. The attributes phase, f, n, NA can be learned.
  \item \texttt{SpatialLightModulator}: Simulates a spatial light modulator (SLM) applied to an incoming field. This element acts as if the SLM is phase only and transmits a field, rather than reflecting it.
  \item \texttt{SeidelAberrations}: Applies Seidel phase polynomial to an incoming field. The coefficients of the polynomials can be learned.
  \item \texttt{ZernikeAberrations}: Applies Zernike aberrations to an incoming field. The coefficients of the polynomials can be learned.
\end{itemize}

% 5. polarizers.py
\paragraph*{5. \textbf{propagation.py}}
\begin{itemize}
  \item \texttt{Propagate}: Free space propagation that can be placed after or between other elements. The method of propagation can be chosen from transform, transfer, exact, asm. The attributes z and n can ba learned
  \item \texttt{KernelPropagate}: This class performs free space propagation using a precomputed propagation kernel (propagator) to propagate an incoming field.
\end{itemize}

% 6. propagation.py
\paragraph*{6. \textbf{samples.py}}
\begin{itemize}
  \item \texttt{ThinSample}: This class is basically a wrapper that wraps the function thin\_sample for training. All the three attributes(absorption, dn, thickness) it takes could be either predefined as an array or a callable for training.
\end{itemize}

% 7. pupils.py
\paragraph*{7. \textbf{sensors.py}}
\begin{itemize}
  \item \texttt{BasicSensor}: The aim of this class is to produce an image representing the intensity of the incoming field with optional shot noise.
\end{itemize}

% 10. sources.py
\paragraph*{8. \textbf{sources.py}}
\begin{itemize}
  \item \texttt{PointSource}: Generates field due to point source a distance ``z`` away. It can also be given pupil. The attributes ``z``, ``n``, ``power``, and ``amplitude`` can be learned.
  \item \texttt{ObjectivePointSource}: Generates field due to a point source defocused by an amount z away from the focal plane, just after passing through a lens with focal length f and numerical aperture NA. The attributes f, n, NA, and power can be learned.
  \item \texttt{PlaneWave}: Generates plane wave of given phase and power.It can also be given pupil and kykx vector to control the angle of the
  plane wave. The attributes kykx, power, and amplitude can be learned. 
  \item \texttt{GenericField}: This is a class built upon the generic\_field function that can be used to generate an electromagnetic field with arbitrary amplitude and phase. Among the attributes, amplitude, phase, and power can be learned.   
\end{itemize}

\paragraph*{8. \textbf{utils.py}}
\begin{itemize}
  \item \texttt{Trainable}: This Trainable class is used to mark and store a value as a trainable parameter which should be learned and optimized.
  \item \texttt{trainable}: This function wraps an input x into a Trainable object to signal to a Chromatix element (lenses, sensors and so on) that x should be used to initialize a trainable parameter.
  \item \texttt{register}: This is used to register attributes of a Flax module, namely the classes defined with nn.Module after the class name, as trainable parameters or fixed state variables.
  \item \texttt{parse\_init}: This function simply defines an inner function init to turn x into an initializer function when x is not a function.
\end{itemize}

\section{Systems Package}
\paragraph*{1. \textbf{amplitude\_masks.py}}
\begin{itemize}
  \item \texttt{AmplitudeMask}: Applies an amplitude mask to an incoming Field by calling the function amplitude\_change from the functional package. The amplitude can be learned pixel by pixel.
\end{itemize}

\section{Performance Considerations}
Chromatix performance depends on array shapes and the complexity of transformations. JAX and XLA (Accelerated Linear Algebra) can compile repeated operations for speed. However, high-resolution fields can be memory-intensive, especially if running multiple parallel wave calculations (e.g., multi-wavelength or batch). GPU usage is recommended.

\chapter{Application Examples}
\section{Example 1: Building Optical System A}
Description, purpose, implementation, and results.

\section{Example 2: Building Optical System B}
Description, purpose, implementation, and results.

\chapter{Conclusion and Future Work}
\section{Conclusion}
Summary of your findings, overall impressions, and effectiveness of Chromatix.

\section{Future Work}
Recommendations for further study or improvements to the library and potential applications.

\chapter{Acknowledgments}
I wish to express my sincere gratitude to my supervisor, Prof. Dr. Vladan Blahnik, for supervising my internship. 

Special thanks to Prof. Dr. , retired professor from [University/Institute], CEO of LightTrans GmbH, for providing the internship topic, invaluable insights, and continuous support throughout the project. 

% Include anyone else you want to thank here.


\bibliographystyle{plain}
\bibliography{references}

\appendix
\chapter{Additional Material}
Include additional scripts, code snippets, or supplementary material here.

\end{document}
